\cotefile{128}
\scanpage{17}
2°) $A'$ \uline{est radiciel sur} $A$, i.e. $\operatorname{Ker}(A'\otimes_{A} A'\to A')$
est nilpotent, donc $\Delta$ \uline{est nilpotent}. Alors
posons
\[
\Omega=\Delta/\Delta^2
\]
et distinguons deux cas

a) $\Omega=0$ d'où $\Delta=\Delta^2$, \deleted{puisque} et $\Delta=0$,
puisque $\Delta$ est nilpotent, donc $A''\simeq A'$, donc
$A^{\unsure{\Delta|1}}\simeq\operatorname{Ker}(A'\rightrightarrows A'')=A'$, donc $\operatorname{long} A'/A=0$.

b) $\Omega\ne0$ ; Notons \unsure{qu'en|que dans} tous cas, \unsure{posons|posant}
pour $x'\in A'$
\[
\begin{aligned}
d(x')&=\delta(x')\bmod\Delta^2\quad\in\Omega\\
&=\varphi(d_{A'/A}(x'))
\end{aligned}
\]
\unsure{l'app.|lorsq.|bien}
où $\varphi:\Omega^1_{A'/A}\to\Omega$ est l'application canonique,
$d$ est une $A$-dérivation de $A'$ dans $\Omega$ :
\[
d(x'y')=x'dy'+y'dx'
\]
d'où
\[
d(\mathfrak m'{}^2)\subset\mathfrak m'\Omega
\]
et comme $\Omega\ne0$, \deleted{\unsure{on} $\mathfrak m'\Omega\ne\Omega$ par Naka} \inserted{un élément de $\mathfrak m'\Omega$ ne}
peut être un générateur de $\Omega$ par Nakayama,
d'où d'après \deleted{$\mathfrak m$} l'hypothèse sur $A'/A$.
\[
\mathfrak m'{}^2\subset A\qquad\text{d'où }\mathfrak m'{}^2\subset\mathfrak m \tag*{$(*)$}
\]
Or comme $A'$ est local à corps résiduels
% archive sheet number: 16
\scanpage{18}
triviale, \deleted{on en conclut} d'où
\[
A'=A+\mathfrak m'
\]
on aura
\[
\mathfrak m A'=\mathfrak m A+\mathfrak m\mathfrak m'\subset\mathfrak m+\mathfrak m'{}^2
\]
et compte tenu de $(*)$ on aura
\[
\mathfrak m A'=\mathfrak m,
\]
C'est à dire on est encore sous les
conditions du N° 1, prop. 3. On conclut
comme dessus (avec les notations du N° 1)
% slanted
\margin{\uline{\unsure{Référence|vérifier} au Lemme préliminaire}}
que la suite
\diagram{\begin{tikzcd}[ampersand replacement=\&,row sep=tiny]
A_0 \arrow[r] \& A'_0 \arrow[r,shift left] \arrow[r,shift right] \& A''_0\\
k \arrow[u,equal] \& \&
\end{tikzcd}}
est encore exacte, et satisfait à \deleted{\unsure{une|encore}} la
\unsure{$\overline{\mathrm m}$|même} condition que la suite du départ :
$x'_0\in A'_0-A_0$ implique que $\delta(x'_0)$ engendre \deleted{$\unsure{\Delta_{A_0}|\Delta}$} $\Delta_0$.

On peut donc encore supposer que $A$ est un corps
$k$, et de plus $A'$ est local à idéal maximal
de carré nul, donc de la forme
\[
A'\simeq D_k(V)\simeq k+V
\]
où $V$ est un vectoriel de dimension finie sur $k$.
% archive sheet number: 17
\scanpage{19}
\deleted{[} Il faut prouver que $V$ est de dim $1$
[il ne peut être de dim $0$, car on aurait
$A'=A$ d'où $\Delta=0$, contrairement à l'hypothèse].

\deleted{Soit $v\in$} Notons que
\[
A'\otimes_k A'\simeq k\,1\otimes1+V\otimes1+1\otimes V+V\otimes V
\]
et que l'idéal engendré par
\[
\delta(v)=v\otimes1-1\otimes v
\]
est \deleted{\unsure{formé}} \inserted{engendré additivement par}
les multiples de $\delta(w)$ et les
$v\otimes w$ et $w\otimes v$.
% archive sheet number: 18
\scanpage{20}
\deleted{\uline{Proposition}} \inserted{\uline{Lemme 1}} Soit $A$ un anneau, $A'$ une $A$-algèbre, $I$ un
idéal \deleted{nilpotent} de carré nul dans $A$, $A_0=A/I$, $A'_0=A'\otimes_A A_0$.
On suppose que, \deleted{$H^1(\unsure{A_0|A'_0}/A_0)=e,\ A_0\xrightarrow{\sim}H^1(A'_0/A_0)$ \unsure{d'une}
$H^1(A'_0/A_0,\unsure{G_a})=0,\ H^1(A'_0/A_0,G_m)$}
posant
\[
\begin{aligned}
B&=H^0(A'/A)=\operatorname{Ker}(A'\rightrightarrows A'')\\
C_0&=H^0(A'_0/A_0)=\operatorname{Ker}(A'_0\rightrightarrows A''_0)
\end{aligned}
\]
on ait
\[
C_0=C_0^*+\operatorname{Im}(B\to C_0)
\]
[N.B. Cette condition est vérifiée si $A$ est noeth. $1/2$ local,
et $A'$ fini sur $A$, car alors $\operatorname{Spec} C_0\to\operatorname{Spec} B_0$ est
radiciel surjectif, donc $B_0$ et $C_0$ sont semi locaux et
l'hom $B_0\to C_0$ induit une \uline{bijection} $\operatorname{Max} C_0\to\operatorname{Max} B_0$,
d'où facilement l'assertion faite]. Sous ces
conditions, on a un isomorphisme canonique
\[
\operatorname{Ker}\big(H^1(A'/A,G_m)\to H^1(A'_0/A_0,G_m)\big)\simeq
\operatorname{Ker}\big(H^1(A'/A,G_a)\to H^1(A'_0/A_0,G_a)\big) \tag*{$(*)$}
\]

\uline{Démonstration}. On a les suites exactes de complexes
\[
0\to I\unsure{C^{\bullet}|C^{\cdot}}(A'/A)\to \deleted{\unsure{K|C}}C^{\bullet}(A'/A)\to \deleted{K}C^{\bullet}(A'_0/A_0)\to0
\]
et
\[
0\to I\deleted{K}C^{\bullet}(A'/A)\to \deleted{K}C^{\bullet}(A'/A;G_m)\to \deleted{K}C^{\bullet}(A'_0/A_0;G_m)\to0
\]
où la deuxième est définie grâce au fait que
$I^2=0$. Par la suite exacte de Cohomologie,
on trouve donc que les deux mbs de $(*)$
% archive sheet number: 19
\scanpage{21}
sont isomorphes respectivement à
\[
\operatorname{Coker}\big(H^0(A'_0/A_0,G_m)\to H^1(IK(A'/A))\big)
\]
et à
\[
\operatorname{Coker}\big(H^0(A'_0/A_0,G_a)\to H^1(IK(A'/A))\big)
\]

Or on a un diagramme commutatif d'applications
\diagram{\begin{tikzcd}[ampersand replacement=\&]
B^* \arrow[r] \arrow[d,"i"] \& C_0^* \arrow[r] \arrow[d,"j"] \& H^1(IK(A'/A)) \arrow[d,"\vert\,\mathrm{id}"] \\
B \arrow[r] \& C_0 \arrow[r] \& H^1(IK(A'/A))
\end{tikzcd}}
où $i,j$ sont les \unsure{immersions|inclusions} canoniques (elles
sont pas multiplicatives, mais cela est sans importance).
Le résultat résulte de cela, et de l'hypothèse
\[
C_0=C_0^*+\operatorname{Im}(B\to C_0).
\]

\uline{Corollaire} 1 Supposons que
\[
H^1(A'_0/A_0,G_a)=0\qquad H^1(A'_0/A_0,G_m)=0
\]
\deleted{\unsure{Pour qu'on ait|Pour qu'il existe}} \inserted{Alors on a} un \deleted{bijection} isomorphisme can.
\[
H^1(A'/A,G_m)\simeq H^1(A'/A,G_a)
\]
donc les relations $H^1(A'/A,G_m)=0$ et
$H^1(A'/A\unsure{,}G_a)=0$ sont équivalentes.
% archive sheet number: 20
\scanpage{22}
\uline{Corollaire} 2 Sous les conditions du lemme\inserted{On suppose que
$\mathbb M_n(C_0)=\mathbb M_n(C_0)^*+\operatorname{Im}(\mathbb M_n(B)\to\mathbb M_n(C_0))$
[ce qui est vérifié sous les conditions signalées au N.B du Lemme 1]}, soit
$n\geq1$ un entier, $\mathbb M_n$ le foncteur ``matrices
carrées de degré $n$ à \unsure{cff. ds|coeff. ds|coeff. dans}'', $\mathrm{Gl}_n=\mathbb M_n^*$, alors on
a une bijection respectant l'unité
\[
\operatorname{Ker}\big(H^1(A'/A,\mathrm{Gl}_n)\to H^1(A'_0/A_0,\mathrm{Gl}_n)\big)\to
\operatorname{Ker}\big(H^1(A'/A,\mathbb M_n)\to H^1(A'_0/A_0,\mathbb M_n)\big)
\]
Par suite, si \deleted{$H^1(A'/A,G_a)=0$, alors}
\inserted{$H^1(A'_0/A_0,\mathrm{Gl}_n)=\{e\}$, $\underbrace{H^1(A'_0/A_0,\mathbb M_n)=\{e\}}_{\substack{\Updownarrow\\H^1(A'_0/A_0)=0}}$,}
on a une bijection
\[
H^1(A'/A,\mathrm{Gl}_n)\xrightarrow{\sim}H^1(A'/A,\mathbb M_n)=\deleted{\unsure{H^1(}}\mathbb M_n(H^1(A'/A))
\]
donc pour que $H^1(A'/A,\mathrm{Gl}_n)=\{e\}$, il f. et s. que
$H^1(A'/A,G_a)=\{e\}$.

Démonstration. Les deux mbs à comparer
sont égaux respectivement à
\[
H^1(I\deleted{\unsure{K}}C(A'/A,\mathbb M_n))/\mathrm{Gl}(n,C_0)
\]
et à
\[
H^1(I\deleted{K}C(A'/A,\mathbb M_n))/\mathbb M_n(C_0)
\]
or les opérations de $\mathrm{Gl}(n,C_0)$ sur $H^1(I\deleted{K}C(A'/A,\mathbb M_n))$
sont induites par celles de $\mathbb M_n(C_0)$. De plus
un élément de $\mathbb M_n(C_0)$ provenant de $\mathbb M_n(B)$ opère
trivialement. Donc il suffit encore de \deleted{\unsure{vérifier que}} \inserted{utiliser l'hypothèse}
% archive sheet number: 21
\scanpage{23}
\uline{Corollaire} 3 Soit $N$ un ens d'indices, \deleted{et supposons}
$E=A^{(N)}$, et supposons que l'on ait
\[
\operatorname{End}_{C_0}(E_{C_0})=\operatorname{Aut}_{C_0}(E_{C_0})
+\operatorname{Im}(\operatorname{End}_{B_0}(E_{B_0})\to\operatorname{End}_{C_0}(E_{C_0}))
\]
Sous ces conditions, on a une bijection can.
respectant les éléments unités
\[
\begin{aligned}
\operatorname{Ker}\big(H^1(A'/A,\underline{\mathrm{Aut}}(E))\to H^1(A'_0/A_0,\underline{\mathrm{Aut}}(E))\big)&\to
\operatorname{Ker}\big(H^1(A'/A,\underline{\mathrm{End}}(E))\to H^1(A'_0/A_0,\underline{\mathrm{End}}(E))\big)\\
&\hspace{5em}\Vert\\
&\big[(\operatorname{Ker}(H^1(A'/A)\to H^1(A'_0/A_0)))^{(N)}\big]^{N}
\end{aligned}
\]

[Donc le noyau du premier mb est trivial [i.e.
les données de descente sur $E_{\!A'}^{\deleted{\unsure{(N)}}}$ qui sont \inserted{telles que sur $E_{A'_0}$}
\inserted{elles soient} effectives et donnent un Module libre \deleted{\unsure{sur le}} $\simeq E_{\!C_0}^{\deleted{\unsure{(N)}}}$
sur $C_0$], sont elles \unsure{$\overline{\mathrm m}\mathrm s$|mêmes} effectives et donnent
un Module libre $\simeq E_B$ sur $B$] ssi
$\operatorname{Ker}(H^1(A'/A)\to H^1(A'_0/A_0))=0$ ] En particulier,
si $H^1(A'_0/\unsure{A'_0|A_0},\underline{\mathrm{Aut}}(E))=\{e\}$ et $H^1(A'_0/A_0)=0$, on a
une bijection
\[
\begin{aligned}
H^1(A'/A,\underline{\mathrm{Aut}}(E))&\simeq H^1(A'_0/A_0,\operatorname{End}(E))\\
&\hspace{3em}\Vert\\
&[H^1(A'/A)^{(N)}]^N
\end{aligned}
\]

Démonstration. $\overline{\mathrm M}$ principe que précédemment,
en utilisant la suite exacte de groupes simpliciaux
\diagram{\scriptsize\begin{tikzcd}[ampersand replacement=\&,column sep=small]
e \arrow[r] \& \deleted{\unsure{K}}C^{\bullet}(A'/A,\underline{\mathrm{Hom}}(E,IE)) \arrow[r]\arrow[d] \& \deleted{\unsure{K}}C^{\bullet}(A'/A,\underline{\mathrm{Aut}}(E)) \arrow[r]\arrow[d] \& \deleted{\unsure{K}}C^{\bullet}(A'_0/A_0,\underline{\mathrm{Aut}}(E)) \arrow[r]\arrow[d] \& 1\\
0 \arrow[r] \& \deleted{\unsure{K}}C^{\bullet}(A'/A,\underline{\mathrm{Hom}}(E,IE)) \arrow[r] \& \deleted{\unsure{K}}C^{\bullet}(A'/A,\underline{\mathrm{End}}(E)) \arrow[r] \& \deleted{\unsure{K}}C^{\bullet}(A'_0/A_0,\underline{\mathrm{End}}(E)) \arrow[r] \& 0
\end{tikzcd}}
% archive sheet number: 22
\scanpage{24}
\uline{Lemme 2} Soit $\deleted{\unsure{A}}B$ un anneau $1/2$ local, $\deleted{\unsure{B}}C$ une algèbre \deleted{finie}
sur $B$, telle que $C$ soit $1/2$ local, et les idéaux
induits par les idéaux maximaux \inserted{\deleted{\unsure{$\mathfrak m$}} $\mathfrak n$} de $C$ sont des
idéaux maximaux \deleted{\unsure{de}} \inserted{$\mathfrak m$} de $B$. Supposons l'une
des conditions suivantes satisfaite

(i) Les corps $k(\mathfrak m)$ sont infinis

(ii) Deux idéaux maximaux \inserted{$\mathfrak n,\mathfrak n'$} distincts de $C$
induisent des idéaux maximaux \inserted{$\mathfrak m,\mathfrak m'$} distincts de $B$,
et pour tt $\mathfrak n$, l'extension résiduelle $k(\mathfrak n)/k(\mathfrak m)$
est radicielle (p. ex. $C$ fini et radiciel sur $B$
$1/2$ local).

Alors pour tt $n\geq1$, et tout élément $u\in\mathbb M_n(C)$,
$\exists\,v\in\mathbb M_n(B)$ tel que $u+\varphi(v)\in\mathbb M_n(C)^*$.

\uline{Dém}
Soient $\mathfrak m_i$ $(1\leq i\leq\unsure{n|m})$ les idéaux maximaux de
$B$ induits par ceux de $C$. Ils sont en nb fini,
donc pour tt\textsuperscript{e} famille d'éléments $v_{\unsure{i|j}}\in\mathbb M_n(k(\mathfrak m_i))$
il existe un $v\in\mathbb M_n(B)$ qui les donne. D'ailleurs,
pour un $w\in\mathbb M_n(C)$, le fait que $w$ soit
inversible ne dépend que de l'inversibilité de $\det w$,
donc que de son image dans $\mathbb M_n(C/\operatorname{rad} C)$! Ces
deux remarques nous ramènent au cas où $C$ est
un composé fini de corps, et $B$ aussi ; On
peut $\overline{\mathrm m}$ supposer $B$ un corps, $B=k$.
% archive sheet number: 23
\scanpage{25}
Soit un $v\in\operatorname{End}_C(E_C)$
Il existe $u\in\operatorname{End}_{\deleted{\unsure{A}}B}(E)$
tel que $v+\deleted{\unsure{u\otimes_A B}}u_C$ soit \uline{inversible}.

\uline{Dém} Si $E\simeq B^{(N)}$, on a $\operatorname{End}_B(E)\simeq E^N$,
donc la remarque sur le relèvt d'éléments $v_i\in
\underline{\mathrm{End}}(E)(k(\mathfrak m_i))$ en un $v\in\underline{\mathrm{End}}(E)(B)$ reste valable
[on est ramené à relever des éléments \deleted{\unsure{des}}
$a_i\in E(k(\mathfrak m_i))$, et finalement, pr $E\simeq B^{(I)}$, des
éléments des $k(\mathfrak m_i)$ eux-$\overline{\mathrm m}\mathrm s$]. De $\overline{\mathrm m}$, grâce à
l'hypothèse que $\operatorname{rad}(C)$ nilpotent, il s'ensuit que
l'inversibilité ne dépend que de $w\bmod\operatorname{rad}(C)$.
Cela nous ramène au cas où $B$, \inserted{$C$} composés de corps,
\deleted{puis finalement}, \inserted{puis} grâce à l'hypothèse (ii), on a
alors $B\xrightarrow{\sim}C$, et on peut encore \unsure{$w|v|u$}$=\mathrm{id}_E-u$.
\deleted{\unsure{au où} $B\unsure{\subset|=|\rightrightarrows}C$}
\deleted{Dans le cas (i), on prend $u=\lambda\,\mathrm{id}_E$, \illegible{1}}

N.B. Il faudrait dégager la possibilité d'une
extension résiduelle \unsure{transcend.t|transcendante} $k(t)$, qui fait
le $\overline{\mathrm m}$ effet que (i) ---
% archive sheet number: 24
\scanpage{26}
\uline{Théorème} $A$ anneau \deleted{artinien}\inserted{$1/2$ local noethérien}, \deleted{\unsure{$A$}} $\mathfrak m$ son radical, $A'$ une algèbre
finie sur $A$. Pour tt entier $n$, on pose $A_n=A/\mathfrak m^{n+1}$, $A'_n=A'/\mathfrak m^{n+1}A'=
A'\otimes_A A_n$. \inserted{a) Soit \unsure{$m|M$} un entier $\geq1$.} Les conditions suivantes sont équivalentes

(i) $\forall n$, on a $H^1(A'_n/A_n)=0$

$(\mathrm{ii}_1)$ $\forall n$, on a $H^1(A'_n/A_n,G_m)=0$
% slanted
\margin{\uline{N.B.} Dans a) les hypothèses peuvent s'affaiblir beaucoup, pour obtenir un énoncé pour \unsure{artinien} uniquement [ $A$ anneau \deleted{$1/2$ local à radical nilpotent}, \inserted{artinien} $A'$ algèbre \uline{quelconque} sur $A$ \ldots ]}

$(\mathrm{ii}_m)$ $\forall n$ on a $H^1(A'_n/A_n,\mathrm{Gl}(m))=\{e\}$

(iii) $\forall n$, désignant par $B^{(n)}$ l'anneau $H^0(A'_n/A_n)$,
le morphisme $\operatorname{Spec} A'_n\to\operatorname{Spec} B^{(n)}$ est un morphisme
de descente effective pour la catégorie fibrée des
Modules \inserted{qu. coh.} plats.

b) Ces conditions impliquent

(i) $H^1(A'/A)=0$

(ii) $H^1(A'/A,G_m)=0$, plus généralement pour tout
entier \unsure{$m|M$}, $H^1(A'/A,\mathrm{Gl}(m))=\{e\}$

(iii) Désignant par $B$ l'anneau $H^0(A'/A)$,
le morphisme $\operatorname{Spec}(A')\to\operatorname{Spec}(B)$ est un
morphisme de descente effective pour la catégorie
des modules plats \tikz[baseline=(t.base)]{\node[draw,ellipse,inner sep=1pt] (t) {\uline{de type fini}};}.
% slanted
\margin{\unsure{est-ce nécessaire}}

c) Supposons $A$ \tikz[baseline=(t.base)]{\node[draw,ellipse,inner sep=1pt] (t) {artinien};}, et les conditions de
a) satisfaites. Alors pour tout $B$-préschéma
% slanted
\margin{\unsure{id}}
$X$ plat sur $\deleted{\unsure{A}}B$, posant $X'=X\otimes_B A'$, le morphisme
$X'\to X$ est un morphisme de descente effective
pour la catégorie fibrée des Modules plats quasi-cohérents.
% archive sheet number: 25
\scanpage{27}
\uline{Démonstration} \inserted{\textcircled{a}} Par les lemmes 1 \inserted{Cor. 2} et 2, l'équivalence des
conditions (i) \deleted{\unsure{(ii)}} $(\mathrm{ii}_m)$ \deleted{\unsure{(iii)}} est immédiate par récurrence
sur $n$, compte tenu que (i) elles sont ttes
vérifiées pour $n=0$ (\uline{descente fid plate}) (ii) la condition
du lemme 1 \deleted{relatif à} $\mathbb M_m(B^{(n)})=\mathbb M_m(B^{(n)})^*+\operatorname{Im}\mathbb M_m(B^{\unsure{(n+1)|(n)}})$
est satisfaite \deleted{pour} $\forall\,n$, en vertu du lemme 2 (ii)
et de la théorie des relations d'équivalence
artiniennes \ldots. La condition $(\mathrm{ii}_m)$ signifie que $\forall n$,
$\operatorname{Spec} A''_n\unsure{-|\to}\operatorname{Spec} B^{(n)}$ est un morphisme de descente
effective \deleted{\unsure{uni}} pour la catégorie des Modules qu. coh.
\uline{libres} de \deleted{type fini} \inserted{rang $m$}. Cela signifie aussi que tte
donnée de descente sur le $A'$-module $A'_n{}^m$,
relative à $A'_n/B^{(n)}$, cette donnée de descente
est effective dans la catégorie des modules \uline{plats} ;
c'est évidemment néc., et c'est suffisant, car
\inserted{je dis que} le module descendu $F_{B^{(n)}}$ sur $B^{(n)}$ (\deleted{\unsure{p.}}plat sur $B^{(n)}$) est
nécessairement libre de rang $m$ \inserted{(car il le devient après chgt de base $B^{(n)}\to A'_n$)}. Comme $B^{(n)}$ est
\margin{lemme à dégager}
artinien, il suffit de voir que le module
réduit mod $\operatorname{rad}(B^{(n)})$ l'est, ce qui nous ramène
au cas où l'anneau de base est \unsure{une} \inserted{somme finie de} corps, puis
un \deleted{somme de} corps, où c'est immédiat \ldots.

Donc il est clair que $(\mathrm{iii})\Rightarrow(\mathrm{ii})_m$. Pour voir que
% archive sheet number: 26
\scanpage{28}
\deleted{\unsure{(ii)|(iii)}}\inserted{(i)} $\Rightarrow(\mathrm{iii})$ \inserted{pour un $n$ donné,}
on note que (i) est stable par changement
plat de la base, et (iii) est \uline{invariant} par
changement \inserted{\uline{fid}} plat de la base. Cela nous ramène au
cas où les extensions résiduelles de $A'/A$ sont ttes
triviales. \deleted{\unsure{Mais alors on}} De plus, on peut évidemment
supposer \inserted{$A=A_n=B^{(n)}$, et} $A$ local [en le remplaçant au besoin
par un localisé]. Alors on voit que si
$E_{A'_n{}^{\deleted{\unsure{m}}}}$ \deleted{plat} sur $A'_n{}^{\deleted{\unsure{m}}}$ est muni d'une donnée de
descente, \deleted{\unsure{alors}} rel à $A'_n{}^{\deleted{\unsure{m}}}/A_n{}^{\deleted{\unsure{m}}}$
[i.e. à $A'_n/B^{(n)}$]
$E_{A'_n}$ est néc. un $A'_n$-module \uline{libre} [car il
suffit de voir que son rang \deleted{\unsure{aux}}\inserted{en les} différents pts
de $\operatorname{Spec} A'_n$ est le $\overline{\mathrm m}$, ce qui nous ramène au
cas où $A$ est le spectre d'un corps, où cela
résulte de l'effectivité de la donnée de descente.]
Je dis qu'on a effectivité pour $A'_n/B^{(n)}$, donc
un module descendu libre de $\overline{\mathrm m}$ rang sur $B^{(n)}$.
Cela résulte du lemme 1, cor. 3, compte tenu
du corollaire au lemme \unsure{?|2.}

On a ainsi prouvé a)

\textcircled{b} Pour prouver (i), comme la formation de
$H^1(A'/A)$ commute à l'extension plate de la base
$A$, on peut supposer $A$ \uline{complet}. Alors \deleted{\unsure{$B|E$}} on a
% archive sheet number: 27
\scanpage{29}
$\unsure{C^{\bullet}|C^{\cdot}}(A'/A)=\varprojlim_n C^{\bullet}(A'_n/A_n)$, donc en vertu de M.L. on
aura
\[
\begin{aligned}
H^1(A'/A)&=H^1(C^{\bullet}(A'/A))=\varprojlim H^1(C^{\bullet}(A'_n/A_n))\\
&=\varprojlim_n H^1(A'_n/A_n)=0
\end{aligned}
\]

On peut aussi déduire formellement (i) de
(ii) sans difficulté (en faisant $m=1$, et $m=2$)

(ii) Va résulter de (iii), par le $\overline{\mathrm m}$ argument que
dans a).

(iii) \deleted{\unsure{So}} On voit comme dans (i) qu'on peut
supposer $A$ \deleted{\uline{local}, \unsure{donc}} \uline{complet}, \deleted{et évidemment
aussi \uline{local}. Soit Enfin, on} On peut aussi supposer
$A=B$, et $A$ \uline{local}. Soit $E'_{\deleted{\unsure{A}}}=E_{A'}$ plat de
type fini sur $A'$, avec donnée de descente, et
soit $E=H^0_D(A'/A,E')$. Il faut prouver
\unsure{au|que} 1°) $E$ plat sur $A$ \qquad 2°) $E\otimes_A A'\to E'$ est un isomorphisme. Or on aura
\[
C^{\bullet}_D(A'/A,E')\simeq\varprojlim_n C^{\bullet}_D(A'_n/A_n,E'_n)
\]
d'où \deleted{\unsure{$E$}} par M.L.
\[
E=H^0_D(A'/A,E')=\varprojlim_n H^0(A'_n/A_n,E'_n)
\]
Posons
\[
E^{(n)}=H^0(A'_n/A_n,E'_n)
\]
alors grâce à \deleted{\unsure{a)}} l'hypothèse a), sous la forme (iii),
on voit que $E^{(n)}$ est un $B^{(n)}$-module plat de
T.f., et $E^{(n)}\otimes_{B^{(n)}} A'_n\to E'_n$ est un isom.
% archive sheet number: 28
\scanpage{30}
Moyennant un lemme facile de passage à la limite,
on en déduit que $E=\varprojlim_n E^{(n)}$ est plat sur
$A=\varprojlim_n B^{(n)}$. Pour vérifier \unsure{enfin|alors} que $E\otimes_A A'\unsure{-|\to} E'$ est
un isomorphisme, on note que les deux mbs sont
resp. les $\varprojlim$ des $E^{(n)}\otimes_{B^{(n)}_{\deleted{\unsure{n}}}} A'_n$ et des $E'_n=E'\otimes_{\unsure{A|A'}} A'_n$,
et l'hom envisagé est la $\varprojlim$ des \uline{isoms}
$E^{(n)}\otimes_{B^{(n)}} A'_n-E'_n$, donc est un \uline{isomorphisme}.

\textcircled{c} \deleted{\unsure{On note|On sait que|Montrons que}} $X'\to X$ est un morphisme de
descente pour la catégorie en question, pour voir que
c'est un morphisme de descente effective, on peut
supposer $X$ affine. De plus, on peut supposer
évidemment $A=B$ [NB, on n'utilise plus les $A_n$ \ldots, mais
l'hypothèse sous la forme
a) (iii) --- ]

\diagram{\begin{tikzcd}[ampersand replacement=\&]
A \arrow[r,dash,very thick] \& A' \arrow[r,dash,very thick,shift left]\arrow[r,dash,very thick,shift right] \& A''
\end{tikzcd}}
% archive sheet number: 29
\scanpage{32}
\[
A'=k+V=k[X_1,\ldots X_n]/(.X_i..)^2
\]

$I$ ``equivalence ideal'' in $A'\otimes_k A'$

Let $F_\alpha(X,Y)=\underset{\text{lin.}}{L'_\alpha(X)}+L''_\alpha(Y)+\underset{\text{bilin.}}{Q_\alpha(X,Y)}$ be a set of generators
\[
p_{13}(I)\subset p_{12}(I)+p_{23}(I)
\]
$\Rightarrow\exists\ G_{\alpha\beta}(X,Y,Z)$ and $H_{\alpha\gamma}(X,Y,Z)$ such that
\[
F_\alpha(X,Z)\equiv\sum_\beta F_\beta(\unsure{x,y|X,Y})
\underset{\substack{\Vert\\\{\ldots+g_{\alpha\beta}(Z)+\ldots\}\\\text{linear}}}{G_{\alpha\beta}(X,Y,Z)}
+\sum_\gamma F_\gamma(Y,Z)
\underset{\substack{\Vert\\\{h_{\alpha\gamma}(X)+\ldots\}\\\text{linear}}}{H_{\alpha\gamma}(X,Y,Z)}
\quad(\mathrm{mod}\,((..X_i.)^2,(.Y)^2,(Z)^{\unsure{2}}))
\]
\[
Q_\alpha(X,Z)=\sum_\beta L'_\beta(X)g_{\alpha\beta}(Z)+\sum_\gamma L''_\gamma(Z)h_{\alpha\gamma}(X)
\]

Take
\[
F'_\alpha(X,Y)=F_\alpha(X,Y)-\sum_\beta F_\beta(X,Y)g_{\alpha\beta}(Y)
-\sum_\gamma F_\gamma(X,Y)h_{\alpha\gamma}(X)=L'_\alpha(X)+L''_\alpha(Y)
\]
This is still set of generators, for:
\[
F_\alpha(X,Y)=F'_\alpha(X,Y)+\sum_\beta F'_\beta(X,Y)g_{\alpha\beta}(Y)+\sum_\gamma F'_\gamma(X,Y)h_{\alpha\gamma}(X)
\]

Next $I\subset\Delta$
\[
\begin{aligned}
F'_\alpha(X,Y)&\equiv\sum_i(X_i-Y_i)\,R_{\alpha i}(X,Y)\qquad
(\mathrm{mod}\,((X_i)^2,(\ldots Y_i)^2))\\
\Rightarrow\quad F'_\alpha(X,Y)&=L'_\alpha(X)-L'_\alpha(Y)
\end{aligned}
\]
% archive sheet number: 31
